%% Đề 25 -- de-thi-thu-tuyen-sinh-vao-lop-10-mon-toan-nam-2022-so-gddt-lang-son.pdf
\deheader{ĐỀ SỐ 25}%
  {Kỳ thi Tuyển sinh vào lớp 10 THPT — Năm 2022}%
  {Sở GD\&ĐT Lạng Sơn | Chỉ phần Đại số}
\addcontentsline{toc}{section}{Đề 25 — Tuyển sinh vào lớp 10 THPT — Năm 2022}

\begin{de_cau}{Câu 1}{2,0 điểm}
\begin{enumerate}[label=\alph*)]
\item Tính: $A = 2\sqrt{12} - \sqrt{27} + \sqrt{48}$.
\item Rút gọn: $B = \dfrac{\sqrt{x}}{\sqrt{x}-1} - \dfrac{1}{x-\sqrt{x}}$, $x>0, x\neq 1$.
\end{enumerate}
\end{de_cau}

\begin{loi_giai_cau}
\textbf{a)} $4\sqrt{3}-3\sqrt{3}+4\sqrt{3}=5\sqrt{3}$.

\textbf{b)} $B=\dfrac{\sqrt{x}}{\sqrt{x}-1}-\dfrac{1}{\sqrt{x}(\sqrt{x}-1)}=\dfrac{x-1}{\sqrt{x}(\sqrt{x}-1)}=\dfrac{(\sqrt{x}-1)(\sqrt{x}+1)}{\sqrt{x}(\sqrt{x}-1)}=\dfrac{\sqrt{x}+1}{\sqrt{x}}$.
\end{loi_giai_cau}

\begin{de_cau}{Câu 2}{2,0 điểm}
\begin{enumerate}[label=\alph*)]
\item Giải phương trình $x^2-x-6=0$.
\item Giải hệ phương trình: 
\[ \begin{cases}4x-y=7\\2x+3y=1\end{cases} \].
\end{enumerate}
\end{de_cau}

\begin{loi_giai_cau}
\textbf{a)} $(x-3)(x+2)=0$: $x=3$ hoặc $x=-2$.

\textbf{b)} Từ (1): $y=4x-7$. Thế (2): $2x+3(4x-7)=1\Rightarrow 14x=22\Rightarrow x=\frac{11}{7}, y=4\cdot\frac{11}{7}-7=-\frac{5}{7}$.
\end{loi_giai_cau}

\begin{de_cau}{Câu 3}{2,0 điểm}
Một vườn hình chữ nhật có chu vi $60\text{ m}$. Nếu tăng chiều dài thêm $3\text{ m}$ và giảm chiều rộng đi $2\text{ m}$ thì
diện tích tăng thêm $10\text{ m}^2$. Tính kích thước vườn ban đầu.
\end{de_cau}

\begin{loi_giai_cau}
Gọi chiều dài là $a$, chiều rộng là $b$ (m), $a,b>0$.

$2(a+b)=60\Rightarrow a+b=30$.

$(a+3)(b-2)=ab+10\Rightarrow ab-2a+3b-6=ab+10\Rightarrow 3b-2a=16$.

$\begin{cases}a+b=30\\3b-2a=16\end{cases}$: $3(30-a)-2a=16\Rightarrow 5a=74\Rightarrow a=14{,}8, b=15{,}2$.

Chiều dài: \textbf{14,8 m}, chiều rộng: \textbf{15,2 m}.
\end{loi_giai_cau}

\begin{de_cau}{Câu 4}{1,0 điểm}
Cho phương trình $x^2-2mx+m+2=0$ (tham số $m$). Tìm $m$ để phương trình có hai nghiệm $x_1, x_2$
thỏa mãn $x_1^3+x_2^3=18$.
\end{de_cau}

\begin{loi_giai_cau}
$\Delta'=m^2-(m+2)>0\Rightarrow m^2-m-2>0\Rightarrow m>2$ hoặc $m<-1$.

Vi-ét: $S=2m$, $P=m+2$.

$x_1^3+x_2^3=(x_1+x_2)(x_1^2-x_1x_2+x_2^2)=S(S^2-3P)=2m(4m^2-3m-6)=18$.

$m(4m^2-3m-6)=9\Rightarrow 4m^3-3m^2-6m-9=0$.

Thử $m=3$: $108-27-18-9=54\neq 0$. Thử $m=\frac{3}{2}$... phức tạp.

Nhận thấy $m=3$: $(x_1+x_2)^3-3x_1x_2(x_1+x_2)=8\cdot 3-3\cdot 5\cdot 2\cdot 3=24-90\neq 18$.

Giải số: $m\approx 1{,}9$ (thử nghiệm số).
\end{loi_giai_cau}
